% =============================================================================
%  2-Desarrollo.tex — Apimetro: arquitectura, pipeline y validación empírica
% =============================================================================

\section{Apimetro: arquitectura e infraestructura tecnológica}

El propósito fundamental de este trabajo es presentar el desarrollo, la arquitectura
y la validación metodológica de \textbf{Apimetro}
\citep{garibaldi2025apimetro}, una infraestructura de datos de software libre
diseñada para resolver la fragmentación de la información de movilidad en la ZMVM.
Apimetro no se concibe únicamente como un servicio de software \textit{back-end},
sino como un motor de democratización de datos geoespaciales cuyo objetivo es
eliminar la barrera técnica del procesamiento de datos espaciales, permitiendo que
investigadores y tomadores de decisiones consuman topologías de red unificadas,
continuas y enriquecidas de manera directa en Sistemas de Información Geográfica
(SIG).

\subsection{Pipeline de datos y stack tecnológico}

El ecosistema opera mediante una arquitectura desacoplada en dos capas.  La primera
corresponde al servicio de unificación de datos espaciales, construido en el
lenguaje de programación Go ---que garantiza alta eficiencia en el procesamiento
concurrente de flujos masivos de datos--- en combinación con una base de datos
relacional espacial PostgreSQL equipada con la extensión PostGIS.

Apimetro ingiere los \textit{feeds} GTFS oficiales de SEMOVI y los somete a un
proceso de normalización y estandarización.  Para resolver los vacíos de información
del Estado de México y corregir trazos geométricos gubernamentales discontinuos
---como el caso del Trolebús Línea~13 o los sistemas de transporte metropolitano
del Estado de México---, la API utiliza de forma estratégica la cartografía
colaborativa de \textbf{OpenStreetMap} como fuente de validación cruzada y datos
complementarios.  Mediante consultas SQL espaciales nativas, Apimetro unifica ambas
fuentes y genera un modelo de datos coherente.

\subsection{Rol de OpenStreetMap en el pipeline}

La contribución de OSM al pipeline de Apimetro opera en tres niveles concretos:

\begin{enumerate}[nosep, leftmargin=1.5em]
  \item \textbf{Geometrías de ruta:}  Cuando los trazos GTFS presentan
    discontinuidades o resolución insuficiente, Apimetro consulta las relaciones
    OSM de tipo \texttt{route=bus}, \texttt{route=subway} y
    \texttt{route=light\_rail} para obtener geometrías continuas de alta
    resolución.  Las relaciones OSM modelan rutas completas como secuencias
    ordenadas de vías (\textit{ways}) y paradas (\textit{nodes}), lo que permite
    reconstruir trazos fidedignos tramo a tramo.
  \item \textbf{Cobertura periférica:}  Los municipios conurbados del Estado de
    México carecen en gran medida de \textit{feeds} GTFS publicados.  Las rutas
    de sistemas como el Mexibús y los colectivos periféricos se obtienen
    íntegramente de OSM, donde mapeadores locales han documentado recorridos que
    las bases oficiales ignoran \citep{haklay2010good}.
  \item \textbf{Validación cruzada:}  Apimetro ejecuta un proceso de
    \textit{conflation} espacial que compara las geometrías GTFS con las de OSM
    mediante distancia de Hausdorff.  Las discrepancias superiores a un umbral
    configurable se marcan para revisión, lo que permite identificar errores
    sistemáticos en las fuentes gubernamentales.
\end{enumerate}

Este enfoque de fuentes duales ---gubernamental y comunitaria--- no solo mejora la
calidad del modelo resultante, sino que devuelve valor a la comunidad OSM al
documentar y validar la infraestructura de transporte con rigor académico.

El sistema expone la red multimodal completa ---Metro, Metrobús, Cablebús, RTP,
Trolebús, Tren Ligero, Tren Suburbano, Mexibús, Colectivo, Microbús y Autobús---
dividida en dos entidades relacionales: \textit{estaciones} (puntos con jerarquía
de transporte y nodos CETRAM identificados) y \textit{líneas} (geometrías
\texttt{MultiLineString} segmentadas tramo a tramo entre estaciones).  Un algoritmo
de \textit{snapping} intermodal con umbral de $\leq 50$~m resuelve las
interconexiones entre sistemas.  La salida se sirve de manera pública a través de
\textit{endpoints} REST en formato GeoJSON nativo desde \texttt{apimetro.dev}.

En total, el motor procesa 22,766 entidades espaciales que se consolidan en un grafo
de \textbf{11,115~nodos} interconectados por \textbf{45,679~aristas}
(25,000 de servicio y 20,000 de transbordo), abarcando \textbf{11~sistemas} de
transporte ---de los cuales el 88\,\% corresponde a transporte de superficie y solo
el 12\,\% a infraestructura masiva de vía exclusiva---, todo en un volumen de
38~MB de datos procesados.  El código es abierto, sin dependencias de software
comercial: soberanía tecnológica completa.

\subsection{Validación empírica: VFTModel como cliente de Apimetro}

Para demostrar la robustez de la infraestructura, se desarrolló el
\textit{VFTModel} (\textit{Vanishing Fig-Tree Model})
\citep{garibaldi2026vft}, un motor analítico en Python que consume directamente
los \textit{endpoints} GeoJSON de Apimetro.  Utilizando librerías de análisis
espacial y redes complejas (NetworkX, GeoPandas, Momepy), el modelo transforma las
geometrías vectoriales en grafos matemáticos ruteables y calcula tres indicadores
clave de eficiencia de red, exportados como capas vectoriales GeoPackage
(\texttt{.gpkg}) directamente visualizables en QGIS
\citep{fortin2016innovative, aldous2019optimal}:

\begin{enumerate}[nosep, leftmargin=1.5em]
  \item \textbf{Cobertura espacial ($C$):} isócronas peatonales de 800~m alrededor
    de cada estación.  Resultado: Milpa Alta presenta solo el 11.75\,\% de cobertura,
    la peor alcaldía de la Ciudad de México.
  \item \textbf{Factor de Desviación ($DF$):} cociente entre la distancia real de
    viaje en la red y la distancia euclidiana.  El promedio de una muestra de
    100~rutas interalcaldías arroja $\bar{DF}=1.545$, lo que significa que cada viaje
    recorre, en promedio, 54\,\% más de lo necesario.  El caso extremo (Milpa Alta
    $\to$ Santa Catarina Yecahuizotl) alcanza $DF=2.52$.
  \item \textbf{Fuerza Capilar ($FC_H$):} centralidad y grado nodal del grafo para
    identificar hubs con mayor concentración de flujos sin alternativa lateral.  El
    hub máximo registra $FC_H=130$ en la zona de Hospitales Sur.
\end{enumerate}

Estos resultados demuestran que la API de Apimetro permite investigación cuantitativa
de nivel académico sin requerir que el investigador programe, limpie archivos de texto
ni repare trazos geométricos rotos.
